\documentclass[10pt,a4paper]{amsart}
\usepackage[margin=19mm]{geometry}
\usepackage{amsmath,amssymb,mathtools}
\usepackage[scr=rsfs,cal=euler]{mathalfa}
\usepackage{xfrac}
\usepackage[unicode,hidelinks,bookmarks=false]{hyperref}
\newcommand{\ie}{i.e.\ }
\newcommand{\cf}{cf.\ }
\newcommand{\resp}{respectively\ }
\newcommand{\Subsection}{\subsection}
\providecommand{\nobreakdash}{\nobreak\hbox{-}}
\setlength{\emergencystretch}{2em}
\multlinegap0pt
\usepackage{amssymb}
\usepackage{mathtools}
\usepackage[shortlabels]{enumitem}
\usepackage{algorithm}\usepackage{algpseudocode}
\usepackage{xcolor}
\newtheorem{theorem}{Theorem}
\newtheorem{assumption}{Assumption}
\DeclareMathOperator{\Cl}{cl}
\newcommand{\E}{\mathbb{E}}
\newcommand{\Pbar}{\bar P_N}
\newcommand{\R}{\mathbb{R}}
\newcommand{\XiSet}{\Xi}
\newcommand{\inner}[2]{\left\langle #1,#2\right\rangle}
\newcommand{\norm}[1]{\left\lVert #1\right\rVert}
\title{Learning the Center and Radius of Wasserstein Ambiguity Sets for Data-Driven Decision Making}
\author{Junjie Guo}
\date{Source: arXiv:2608.18123v1 (CC BY 4.0)}
\begin{document}
\maketitle

\noindent\emph{Source and licence.} Learning the Center and Radius of Wasserstein Ambiguity Sets for Data-Driven Decision Making. Version arXiv:2608.18123v1 (CC BY 4.0), distributed by arXiv under the Creative Commons Attribution 4.0 International licence, http://creativecommons.org/licenses/by/4.0/, as recorded in the arXiv licence metadata for this exact version and shown on the abstract page https://arxiv.org/abs/2608.18123v1. Full licence notice: \texttt{licenses/arxiv-2608.18123-CC-BY-4.0.txt}.\par\smallskip

\noindent\textit{Editorial scope.} Counted unit: the article's theorem thm:convex-reform-general (source0.tex lines 291--303). The article's complete original proof of it is delivered with it (source0.tex lines 305--334, one \texttt{proof} environment of 30 lines). No mathematics is written, repaired or re-derived: the statement and the proof are the article's own lines, in the article's own notation, and no retained line is substituted or re-flowed.  Retained uncounted as the source-local context the proof needs: lines 238--240; lines 242--248; lines 244--247; lines 283--289.  Nothing was omitted from any retained span, so the package carries no omission record.  No retained line contains a citation, so the wrapper carries no bibliography and no bibliography entry.  No retained line contains a figure environment, an image-inclusion command or drawing code, so no image is delivered and the package declares no asset dependency.  Editorial formatting only, in addition: an AMS \texttt{amsart} preamble that restates the article's own declarations and macros used by the retained lines, namely the environments \texttt{theorem}, \texttt{assumption} on separate counters, together with the standard packages those lines need.  The retained environments print in this document as Theorem 1, Assumption 1 in order of appearance, whereas the article prints the same items with the numbering of its own full document.  The complete native source of the article as distributed in the arXiv e-print for this exact version is bundled unchanged as \texttt{sources/207968/source0.tex}.
\begin{equation}\label{eq:wce-general}
    \Phi_{\varepsilon}(\ell;\Pbar)=\sup_{Q\in B_\varepsilon(\Pbar)}\E_Q[\ell(\xi)].
\end{equation}

\begin{theorem}[Dual worst-case expectation]\label{thm:dual-general}
Suppose strong duality holds for the moment problem in \eqref{eq:wce-general}; in particular, this is true under the standard upper-semicontinuity and linear-growth conditions used in Wasserstein DRO. Then
\begin{equation}\label{eq:dual-general}
    \Phi_{\varepsilon}(\ell;\Pbar)
    =\inf_{\lambda\ge0}\left\{\lambda\varepsilon+\sum_{j=1}^Mp_j\sup_{\xi\in\XiSet}\left(\ell(\xi)-\lambda\norm{\xi-\zeta_j}\right)\right\}.
\end{equation}
\end{theorem}

\begin{equation}
    \Phi_{\varepsilon}(\ell;\Pbar)
    =\inf_{\lambda\ge0}\left\{\lambda\varepsilon+\sum_{j=1}^Mp_j\sup_{\xi\in\XiSet}\left(\ell(\xi)-\lambda\norm{\xi-\zeta_j}\right)\right\}.
\end{equation}

\begin{assumption}[Concave-piece loss]\label{ass:concave-piece}
The uncertainty set $\XiSet\subseteq\R^m$ is closed and convex. The loss admits the representation
\begin{equation}\label{eq:max-concave}
    \ell(\xi)=\max_{k\le K}\ell_k(\xi),
\end{equation}
where $-\ell_k$ is proper, closed, and convex for every $k$.
\end{assumption}

\begin{theorem}[Finite convex reformulation]\label{thm:convex-reform-general}
Under Assumption \ref{ass:concave-piece} and the strong-duality conditions of Theorem \ref{thm:dual-general},
\begin{equation}\label{eq:convex-reform-general}
\Phi_{\varepsilon}(\ell;\Pbar)=
\begin{aligned}[t]
\inf_{\lambda,s_j,z_{jk},\nu_{jk}}\quad
&\lambda\varepsilon+\sum_{j=1}^Mp_js_j\\
\mathrm{s.t.}\quad
&[-\ell_k]^*(z_{jk}-\nu_{jk})+\sigma_{\XiSet}(\nu_{jk})-\inner{z_{jk}}{\zeta_j}\le s_j,\quad \forall j,k,\\
&\norm{z_{jk}}_*\le\lambda,\quad \forall j,k.
\end{aligned}
\end{equation}
\end{theorem}

\begin{proof}
Starting from \eqref{eq:dual-general}, introduce epigraph variables $s_j$:
\begin{equation}
\inf_{\lambda\ge0,s_j}\left\{\lambda\varepsilon+\sum_jp_js_j:
\sup_{\xi\in\XiSet}\left(\ell(\xi)-\lambda\norm{\xi-\zeta_j}\right)\le s_j,\ \forall j\right\}.
\end{equation}
Using $\ell=\max_k\ell_k$, the constraints are equivalent to
\begin{equation}\label{eq:semiinf}
    \sup_{\xi\in\XiSet}\left(\ell_k(\xi)-\lambda\norm{\xi-\zeta_j}\right)\le s_j,
    \qquad \forall j,k.
\end{equation}
By the definition of the dual norm,
\begin{equation}
    \lambda\norm{\xi-\zeta_j}=\sup_{\norm{z}_{*}\le\lambda}\inner{z}{\xi-\zeta_j}.
\end{equation}
Therefore, under the concavity/convexity conditions and the minimax theorem,
\begin{align*}
\sup_{\xi\in\XiSet}\left(\ell_k(\xi)-\lambda\norm{\xi-\zeta_j}\right)
&=\inf_{\norm{z_{jk}}_*
\le\lambda}\sup_{\xi\in\XiSet}\left(\ell_k(\xi)-\inner{z_{jk}}{\xi-\zeta_j}\right)\\
&=\inf_{\norm{z_{jk}}_*
\le\lambda}\left\{[-\ell_k+\chi_{\XiSet}]^*(-z_{jk})+\inner{z_{jk}}{\zeta_j}\right\}.
\end{align*}
Replacing $z_{jk}$ by $-z_{jk}$ and using inf-convolution of conjugates,
\begin{equation}
[-\ell_k+\chi_{\XiSet}]^*(z_{jk})
=\Cl\inf_{\nu_{jk}}\left\{[-\ell_k]^*(z_{jk}-\nu_{jk})+\sigma_{\XiSet}(\nu_{jk})\right\}.
\end{equation}
Substituting this expression into the epigraph constraints yields \eqref{eq:convex-reform-general}. The closure does not change the resulting inequality constraints under standard lower-semicontinuity conventions.
\end{proof}

\end{document}
